\chapter{Conclusions and outlook}

\section{Summary}

The main results of these notes fit into five statements. Each one is listed below together with the chapter
where it was built, so any one of them can be traced back to its full derivation.

\begin{itemize}
\item In the $G_N\to0$ limit, a general bulk subregion is defined by a boundary operator algebra —
typically one with no direct boundary geometric description at all — carrying an emergent type
$\mathrm{III}_1$ structure, sourced by the infinite long-range entanglement that only appears in the strict
large-$N$ limit (Chapters~6 and~7). This machinery is not restricted to the strict Einstein
($\lambda\to\infty$) limit; it extends, with modification, into the stringy regime (Chapter~8).
\item Bulk causal structure is encoded in boundary \emph{commutant} structure — not merely boundary
causality (which only forces spacelike-separated operators to commute), but the richer, timelike commutant
structure that subregion-subalgebra duality inherits from the bulk. Quantitative tools like the causal depth
parameter $T(t)$ (Chapter~8) make this a checkable, numerical diagnostic of horizons and global causal
structure, and these tools too extend into the stringy regime.
\item Modular flow, and its refinement half-sided modular flow, characterize the emergence of time itself in
the bulk — resolving the meeting-behind-the-horizon puzzle via emergent Kruskal time, and motivating a
sharpened, algebraic version of the ER$=$EPR proposal (both in Chapter~8).
\item Large-$N$ boundary algebras split cleanly into two kinds. An \textbf{entanglement wedge algebra} $X$ is,
by definition, the large-$N$ limit of a genuine finite-$N$ algebra $B$ — it always admits a finite-$N$
ancestor. A \textbf{causal wedge algebra} $Y$, by contrast, is intrinsically a large-$N$, semiclassical
construct, built directly from single-trace operators via the extrapolate dictionary, with no finite-$N$
definition of its own. Neither is guaranteed to have a bulk geometric meaning once you leave the strict
Einstein regime — and both constructions apply just as well to non-gravitational systems entangled with a
holographic partner (Chapter~7), with no need for the non-gravitational side to have its own large-$N$
parameter at all.
\item Simple operator-algebraic toy models — static observers, dynamical observers, and exactly solvable JT
gravity (all in Chapter~9) — turn this machinery into concrete physics: the physical origin of generalized
gravitational entropy for black holes, de~Sitter space, and general regions; the emergence of a cosmological
horizon as an observer's mass is taken to infinity; and, in JT gravity, a complete, finite-$G_N$ example where
the operator algebra of quantum gravity is known in exact closed form, even though the Hilbert space itself
does not factorize.
\end{itemize}

\section{The mathematical structure of quantum gravity}

This final section is more speculative than the rest of the notes. It returns to the question raised in
Chapter~1 — \emph{what happens to the wavefunction, and to the Hilbert space itself, once you take the algebra
seriously as the fundamental object?} — and answers it at the largest scale available.

\subsection*{Why quantum gravity cannot have one single Hilbert space}

Ordinary quantum mechanics treats the Hilbert space as fundamental: states are density operators living in it,
observables are operators acting on it. There is a concrete reason why this cannot be the right starting point
for quantum gravity. It is widely believed that no physical process in quantum gravity can change a
spacetime's \emph{asymptotic} structure — there is no physical process that takes a state of type IIB string
theory on $\text{AdS}_5\times S^5$ to a state on $\text{AdS}_3\times S^3\times K3$, or to either of these from
ten-dimensional flat Minkowski space. Since a genuine physical process is exactly what it means for two states
to live in the same Hilbert space (you can always in principle evolve from one to the other), \textbf{each of
these asymptotic backgrounds must correspond to a genuinely separate Hilbert space}. This is not a conjecture;
it is already visible directly in the existing dictionary: AdS$_5\times S^5$ is dual to a four-dimensional CFT
with its own Hilbert space, AdS$_3\times S^3\times K3$ to a two-dimensional CFT with a completely different
one, and nothing maps a state of one to a state of the other. Yet all of these are supposed to be different
vacua of the very same underlying theory — type IIB string theory. \textbf{Without one single, global Hilbert
space to hold all of these together, what mathematical object is left to call ``the theory''?}

\subsection*{The resolution you have already met, twice}

This is exactly the same shape of question that opened Chapter~1, and the resolution is precisely the demotion
of the Hilbert space worked through in detail in Chapter~2, applied now at the grandest possible scale. Look
first at a much smaller-scale version of the same puzzle, already familiar from ordinary QFT in curved
spacetime: a free scalar field in flat Minkowski space can be quantized using ordinary Minkowski time,
\emph{or} using Rindler time (Chapter~4). Both are perfectly sensible choices, describing different physical
situations (an inertial versus an accelerated observer), and yet they produce genuinely, unitarily
\emph{inequivalent} Hilbert spaces. You could declare these to be two different theories, but that is clearly
the wrong instinct — there is an obvious sense in which it is the same free scalar field theory, just quantized
around two different reference states. \textbf{Algebraic QFT resolves this by defining the theory as a pair
$(\Alg,\Sscr)$ — an algebra $\Alg$ of field operators, together with a space $\Sscr$ of allowed states (in the
abstract, linear-functional sense of Chapter~2, not vectors in any particular Hilbert space) — and letting the
Hilbert space be a \emph{derived}, state-dependent object, produced fresh by the GNS construction of Chapter~2
for whichever state $\omega\in\Sscr$ you happen to pick.} Minkowski quantization and Rindler quantization are
simply two different GNS representations of the very same underlying algebra $\Alg$ — genuinely different
Hilbert spaces, but manifestly the same theory, because the algebra never changed.

\textbf{The proposal is to do exactly this, one level up, for the whole of quantum gravity.} Postulate that a
quantum-gravity theory at finite $G_N$ is specified by a single abstract $*$-algebra $\Alg$ (not tied to any
particular asymptotic structure) together with some allowed collection of states on it. Different asymptotic
structures — AdS$_5\times S^5$, AdS$_3\times S^3\times K3$, flat space, de~Sitter — simply correspond to
\emph{different states} $\omega$ on this one algebra, each producing its own GNS Hilbert space $\HH_\omega$ via
exactly the construction of Chapter~2, with the corresponding boundary CFT (where one is known) simply
\emph{being} that GNS representation. Concretely: $\omega_{\text{AdS}_5\times S^5}$ is the state whose GNS
Hilbert space is that of $\mathcal N=4$ super-Yang-Mills; $\omega_{\text{AdS}_3\times S^3\times K3}$ is the
state whose GNS space is the dual two-dimensional CFT's; and, in principle, flat space and de~Sitter correspond
to further states of the same underlying algebra $A_{\text{IIB}}$, even though an explicit boundary description
for those cases is not yet known.

\begin{physconnect}[: levels of algebraic emergence]
This is the same logical move made twice already in these notes. Lined up side by side, the three levels are
each a strictly more drastic version of the same idea:
\begin{itemize}
\item \textbf{States (Chapter~2, Slavnov's statistical picture)}: an ordinary quantum \emph{state} $\omega$
(or $\rho$, or $\ket\psi$) is not fundamental — it is a derived bookkeeping device (an ensemble average, in
Slavnov's language) built on top of the more primitive pairing of algebra and individual measurement outcome.
\item \textbf{Hilbert spaces (Chapter~2's GNS construction, and the Minkowski--Rindler example just above)}:
the \emph{Hilbert space itself} is not fundamental — a single algebra $\Alg$ can produce many inequivalent
Hilbert spaces, one per choice of reference state, via GNS.
\item \textbf{Spacetime backgrounds (this section)}: even the question ``which spacetime background am I in''
is not fundamental data fed into the theory — it, too, is just a choice of state on one underlying algebra,
with the entire geometric structure of a spacetime (AdS, flat space, de~Sitter, and everything built on top of
it in these notes — causal structure, horizons, entropy) emerging as a \emph{consequence} of that choice, not an
ingredient of it.
\end{itemize}
Nothing about the mathematics changed between these three instances — it is the identical GNS construction,
applied at the level of an individual measurement, then at the level of an ordinary quantum system, then at
the level of an entire spacetime background. What changed is only how much structure you are willing to
regard as emergent rather than fundamental, and each chapter of these notes have pushed that boundary a little
further out.
\end{physconnect}

At present, there is no known background-independent definition of the algebra $A_{\text{IIB}}$ itself, nor a
full characterization of which states are allowed. String field theory is a plausible route toward one, since
for any fixed background state it already supplies both the Hilbert space and the background-dependent
algebra, with its equations of motion available to help identify further consistent states systematically.
The algebraic program developed in these notes — subregion-subalgebra duality, modular time, the crossed
product, generalized entropy from observer dressing — can be read as a first, concrete step toward exactly
this larger goal: a formulation of quantum gravity in which asymptotic structure, the cosmological constant,
the Hilbert space, and the very number of degrees of freedom are not fixed ingredients handed to the theory in
advance, but \emph{derived, state-dependent consequences} of a single, more primitive algebraic structure.

\section*{Further reading}

Two technical computations are not reproduced in these notes; both can be found in~\cite{Liu2025}. The first
consists of further finite-dimensional examples of infinite entanglement coexisting with a factorizable Hilbert
space — fine-tuned counterexamples to the general expectation, discussed in Chapter~2, that infinite
entanglement generically prevents factorization. The second is the explicit solution of the KMS relation that
gives $\widehat\Delta=\Delta_\Psi$ for the crossed-product algebra, the result used without derivation in
Chapter~5. Neither changes the conceptual picture developed here.

\bigskip
\noindent These notes began with a wavefunction and asked what it really is. The answer that emerged, one
worked example at a time, is that the wavefunction, the Hilbert space it lives in, and finally the spacetime
itself are each the most convenient representation of something more primitive: an algebra of observables,
together with a choice of state on it. Which parts of the familiar picture are fundamental, and which are
representations of something deeper, is the same question at every level — for a single spin, for a quantum
field in a Rindler wedge, for a black hole, and for the universe as a whole.
